\section{A genuine C-to-S bridge while the opposite deficit is rising}
This section uses exact source coefficients and their monotonicity,
without any forward cumulative log-concavity assumption.

\begin{theorem}[Nonrecovering reverse deficit from one-arm C]
Let $A=F_Z$, $B=F_Y$ be actual increasing nonrectangular arms, and
assume $A$ has property C. Set $D_i=(A_i-A_{i-1})^+$,
$D_0=D_1=0$, $E_B=1-z+(2z-1)B$, and $H=E_BA$.
At a forward comparison $j\ge2$, assume
$D_jD_{j+1}>0$ and $E_B$ is nondecreasing throughout ranks
zero through $j+1$. If
\[
                    2(AB)_j-(AB)_{j+1}\ge0,
\]
then
\[
            2(AB)_{j-1}-(AB)_j
                       \le2(AB)_j-(AB)_{j+1}.           \tag{1}
\]
In particular the deficit monotonicity premise holds whenever $B$
is still strictly rising at $j+1$, by the established single-arm
ordinary/deficit mode ordering. If both ordinary factors are still
strictly rising, the theorem applies whenever either one has C.
\end{theorem}
\begin{proof}
Write $w_r=E_B(r)-E_B(r-1)$. Nonrectangularity gives $w_r=0$
for $r<3$; the prefix hypothesis gives $w_r\ge0$ through $j+1$.
The exact source identities are
\[
 H_j=\sum_{r=3}^j w_rG_A(j-r),\qquad
 H_{j+1}=\sum_{r=3}^{j+1}w_rG_A(j+1-r).
\]
Multiply the shift-$r$ C inequalities by $w_r$ and sum, retaining
the nonnegative newly born $r=j+1$ term. This proves
$H_{j+1}D_j\ge H_jD_{j+1}$. Also
$H_{j+1}-H_j=\sum_rw_rA_{j+1-r}\ge0$.

The exact identity
\[
                    (2z-1)AB=E_BA-(1-z)A
\]
identifies the two deficits in (1) with $H_j-D_j$ and
$H_{j+1}-D_{j+1}$. The latter is nonnegative by hypothesis.
These facts imply its monotonicity by the positive decomposition
\[
\begin{split}
 H_{j+1}\bigl((H_{j+1}-D_{j+1})-(H_j-D_j)\bigr)
  ={}&(H_{j+1}D_j-H_jD_{j+1})\\
    &+(H_{j+1}-H_j)(H_{j+1}-D_{j+1})\ge0.
\end{split}
\]
Here $H_{j+1}\ge D_{j+1}>0$, so division proves (1).
This step uses BOTH the unclipped ratio comparison and monotonicity
of the actual source convolution; clipping alone would not suffice.
\end{proof}

\begin{corollary}[Correct coupled S on the two-rising-factor branch]
Put $p=q_Zq_Y$, $a=\operatorname{wt}Z$, $b=\operatorname{wt}Y$,
and $R_i=2p_i-p_{i-1}$. At a usual S crossing index $n$ set
$j=a+b-n$. If $j\ge2$, both $F_Z,F_Y$ are still strictly rising
at $j+1$, at least one of them has C, and $R_n\ge0$, then
\[
                         R_{n+1}\le R_n.                \tag{2}
\]
Assume, as in the A2 induction, that the OLD weighted positive
differences $2^iu_i$ are unimodal. For $k<\ell$, with $\ell$ their
first mode, one has $u_k\le2u_{k+1}$. Formula (2) is exactly
$v_n\le2v_{n+1}$, where $v_n=p_{n-1}-p_n$ in the crossing
range. Adding gives
\[
                    u_k+v_n\le2(u_{k+1}+v_{n+1}),
\]
which is S. Strict weighted fall is not needed inside this positive
branch, although it cannot be removed from general S.
\end{corollary}

The completed C theorem through five non-one parts makes the C premise
UNCONDITIONAL whenever either old arm has at most five such parts;
the other arm is unrestricted. The OLD A2 premise remains an inductive
input: the corollary does not assert A2 for arbitrary products or the
original conjecture. General C alone would close this two-rising-factor
branch at every size, not the remaining branches where an individual
factor has already ceased rising.

The companion \texttt{zero\_head\_C\_bridge.tex} combines this with the
local first-defect branch and a double-defect convolution bound. It
closes the ENTIRE correctly conditioned zero-head S range whenever
both old arms have at least two non-one parts and one has C, without
requiring OLD A2 in that zero-head case.

\texttt{CToCoupledSRisingWindow.lean} proves the complete arbitrary-real
monotone-excess implication by the displayed positive decomposition and
the S sum implication, and the integer double-defect/positive-slopes
implication. The unbounded weighted source sum is the existing
\texttt{weighted\_ratio\_sum} certificate from
\texttt{ClippedForwardRatio.lean}. Exact composition source identities,
the known deficit mode ordering and reversal give the written assembly.
No end-to-end composition formalization or general A2/UI solution is
claimed. This is a conditional bridge into an actual weighted-difference
resource, rather than an identification of cumulative LC with coefficient UI.
